
La interacción con células neuronales es una tarea que requiere restricciones temporales muy estrictas, 
por lo que para su estudio, se deben emplear sistemas y herramientas que permitan enviar y registrar 
señales en tiempo real, con precisión en la escala de los milisegundos, además de realizar los cálculos necesarios para implementar interacciones realistas con las neuronas vivas. 
Esta necesidad va en incremento, dada su implicación en múltiples aplicaciones. En estas, también se 
necesitan realizar otras tareas relacionadas con el análisis de la dinámica neuronal y el objetivo de los experimentos, por lo que no se puede 
dejar espacio al sistema operativo para que interrumpa los procesos encargados de realizarlas.

Por tanto, es necesario investigar las distintas opciones y posibilidades 
que van surgiendo en relación a dichos sistemas de tiempo real, además de aquellas herramientas que faciliten las tareas 
anteriormente mencionadas y el mantenimiento a lo largo del tiempo de las mismas.
